\documentclass{standalone}
\usepackage{circuitikz}
\usetikzlibrary{calc}
\definecolor{circuitnotemark}{HTML}{FE00FE}
\begin{document}
\begin{circuitikz}[american, line width=0.8pt]
\ctikzset{bipoles/length=1.2cm}
\foreach \x in {0,2,4} {\foreach \y in {0,-2,-4,-6,-8,-10} {\fill[gray, opacity=0.35] (\x,\y) circle (1.1pt);}}
\node[gray, font=\scriptsize] at (0,0.84) {1};
\node[gray, font=\scriptsize] at (2,0.84) {2};
\node[gray, font=\scriptsize] at (4,0.84) {3};
\node[gray, font=\scriptsize, anchor=east] at (-0.84,0) {a};
\node[gray, font=\scriptsize, anchor=east] at (-0.84,-2) {b};
\node[gray, font=\scriptsize, anchor=east] at (-0.84,-4) {c};
\node[gray, font=\scriptsize, anchor=east] at (-0.84,-6) {d};
\node[gray, font=\scriptsize, anchor=east] at (-0.84,-8) {e};
\node[gray, font=\scriptsize, anchor=east] at (-0.84,-10) {f};
\coordinate (a1) at (0,0);
\coordinate (a3) at (4,0);
\draw (a1) to[R, l_=$R_{1}$, a^=$10\,\mathrm{k}\Omega$] (a3); % line 3
\path (0,-2.395) rectangle (1.778,-1.605); % line 5
\node[anchor=west, inner sep=0, circuitnotemark, font=\footnotesize] at (0,-2) {X}; % line 5
\path (0,-4.494) rectangle (2.064,-3.506); % line 6
\node[anchor=west, inner sep=0, circuitnotemark, font=\normalsize] at (0,-4) {X}; % line 6
\path (0,-6.593) rectangle (3.556,-5.407); % line 7
\node[anchor=west, inner sep=0, circuitnotemark, font=\large] at (0,-6) {X}; % line 7
\path (0,-8.853) rectangle (3.566,-7.147); % line 8
\node[anchor=west, inner sep=0, circuitnotemark, font=\LARGE] at (0,-8) {X}; % line 8
\path (0,-11.229) rectangle (5.53,-8.771); % line 9
\node[anchor=west, inner sep=0, circuitnotemark, font=\Huge] at (0,-10) {X}; % line 9
\path (6,-7.408) rectangle (13.531,0.593); % line 10
\node[anchor=west, inner sep=0, circuitnotemark, font=\large] at (6,0) {X}; % line 10
\node[anchor=west, inner sep=0, circuitnotemark, font=\large] at (6,-0.487) {X}; % line 10
\node[anchor=west, inner sep=0, circuitnotemark, font=\large] at (6,-0.974) {X}; % line 10
\node[anchor=west, inner sep=0, circuitnotemark, font=\large] at (6,-1.46) {X}; % line 10
\node[anchor=west, inner sep=0, circuitnotemark, font=\large] at (6,-1.947) {X}; % line 10
\node[anchor=west, inner sep=0, circuitnotemark, font=\large] at (6,-2.434) {X}; % line 10
\node[anchor=west, inner sep=0, circuitnotemark, font=\large] at (6,-2.921) {X}; % line 10
\node[anchor=west, inner sep=0, circuitnotemark, font=\large] at (6,-3.408) {X}; % line 10
\node[anchor=west, inner sep=0, circuitnotemark, font=\large] at (6,-3.895) {X}; % line 10
\node[anchor=west, inner sep=0, circuitnotemark, font=\large] at (6,-4.381) {X}; % line 10
\node[anchor=west, inner sep=0, circuitnotemark, font=\large] at (6,-4.868) {X}; % line 10
\node[anchor=west, inner sep=0, circuitnotemark, font=\large] at (6,-5.355) {X}; % line 10
\node[anchor=west, inner sep=0, circuitnotemark, font=\large] at (6,-5.842) {X}; % line 10
\node[anchor=west, inner sep=0, circuitnotemark, font=\large] at (6,-6.329) {X}; % line 10
\node[anchor=west, inner sep=0, circuitnotemark, font=\large] at (6,-6.816) {X}; % line 10
\node[anchor=south west, inner sep=2pt, circuitnotemark, font=\LARGE\bfseries] (circuittitle) at (current bounding box.north west) {X};
\path (circuittitle.south west) rectangle ++(5.578,0.853);
\end{circuitikz}
\end{document}
